\documentclass{article}
\usepackage{amsmath, amssymb}
\usepackage{booktabs}
\usepackage{graphicx}
\usepackage{hyperref}
\usepackage{xcolor}

\title{PyDESeq2 Aberrant LFC Shrinkage: A Case Study}
\author{crispyx Development Team}
\date{January 2026}

\begin{document}

\maketitle

\begin{abstract}
We document a significant numerical issue in PyDESeq2's apeGLM log-fold-change (LFC) shrinkage implementation that can cause LFC magnitudes to \emph{increase} rather than decrease after shrinkage. This aberrant behavior affects up to 21\% of genes in single-cell CRISPR screen data when the dispersion trend fitting fails. We provide specific examples and compare against crispyx's more robust implementation.
\end{abstract}

\section{Introduction}

The apeGLM (approximate posterior estimation for GLM) method \cite{apeglm} is designed to shrink log-fold-change estimates toward zero, reducing noise in differential expression results while preserving biologically significant effects. The key property of proper shrinkage is:
\begin{equation}
    |\hat{\beta}_{\text{shrunk}}| \leq |\hat{\beta}_{\text{MLE}}|
\end{equation}
where $\hat{\beta}_{\text{MLE}}$ is the maximum likelihood estimate and $\hat{\beta}_{\text{shrunk}}$ is the posterior mode under a Cauchy prior.

We observed that PyDESeq2's implementation can violate this property, producing ``anti-shrinkage'' where $|\hat{\beta}_{\text{shrunk}}| > |\hat{\beta}_{\text{MLE}}|$.

\section{Root Cause: Dispersion Trend Fitting Failure}

When PyDESeq2's dispersion trend fitting fails to converge, it falls back to a mean-based dispersion trend. This can lead to poorly estimated gene-wise dispersions, which in turn causes the apeGLM optimization to converge to incorrect values.

The warning message from PyDESeq2:
\begin{verbatim}
UserWarning: The dispersion trend curve fitting did not converge. 
Switching to a mean-based dispersion trend.
\end{verbatim}

This warning is often silently ignored, but has severe downstream consequences for LFC shrinkage.

\section{Quantitative Analysis: Adamson Subset Dataset}

We analyzed the Adamson subset dataset, a single-cell CRISPR screen dataset with 1,716 cells and 11,630 genes.

\subsection{Summary Statistics}

\begin{table}[h]
\centering
\caption{Shrinkage aberration rates for different perturbations}
\label{tab:summary}
\begin{tabular}{lccccc}
\toprule
Perturbation & Total Genes\textsuperscript{a} & $|$shrunk$|$ $>$ $|$raw$|$ & $>2\times$ inflation & $>5\times$ inflation & Max ratio \\
\midrule
AARS vs control  & 5,527 & 1,186 (21.5\%) & 182 & 16 & 13.74$\times$ \\
AMIGO3 vs control & 2,367 & 19 (0.8\%) & 0 & 0 & 1.62$\times$ \\
\bottomrule
\multicolumn{6}{l}{\textsuperscript{a}Genes with $|$raw LFC$|$ $>$ 0.1}
\end{tabular}
\end{table}

The AARS perturbation shows severe aberrant shrinkage, with 21.5\% of genes having inflated LFC after shrinkage. Some genes show up to 13.74$\times$ inflation.

\subsection{Specific Examples: AARS vs Control}

\begin{table}[h]
\centering
\caption{Examples of PyDESeq2 aberrant shrinkage (AARS vs control)}
\label{tab:examples_aars}
\begin{tabular}{lcccccc}
\toprule
Gene & Raw LFC & Raw SE & Shrunk LFC & Shrunk SE & Inflation & Dispersion \\
\midrule
AKNA     & 0.299 & 0.707 & 4.111 & 1.201 & 13.74$\times$ & 53.12 \\
MILR1    & 0.294 & 0.661 & 2.978 & 0.947 & 10.12$\times$ & 45.86 \\
NDUFA4L2 & 0.276 & 0.670 & 2.631 & 0.933 & 9.53$\times$  & 47.19 \\
APOL6    & 0.252 & 0.688 & 2.248 & 0.949 & 8.93$\times$  & 50.03 \\
HRK      & 0.302 & 0.621 & 2.450 & 0.839 & 8.12$\times$  & 39.99 \\
CNN1     & $-$0.274 & 0.574 & $-$2.208 & 0.773 & 8.06$\times$ & 33.67 \\
TMEM154  & 0.431 & 0.654 & 3.441 & 0.898 & 7.99$\times$  & 44.14 \\
MYO1F    & 0.444 & 0.637 & 3.288 & 0.858 & 7.41$\times$  & 41.66 \\
FAM129A  & 0.567 & 0.670 & 4.115 & 0.937 & 7.26$\times$  & 46.20 \\
LGMN     & $-$0.255 & 0.454 & $-$1.476 & 0.567 & 5.78$\times$ & 19.61 \\
\bottomrule
\end{tabular}
\end{table}

\textbf{Key observations:}
\begin{itemize}
    \item All aberrant genes have \emph{high dispersions} ($\alpha > 30$), indicating noisy expression patterns
    \item The raw LFC values are relatively small ($|$LFC$|$ $<$ 0.6), with large standard errors
    \item After ``shrinkage'', the LFC magnitudes increase dramatically (up to 13.74$\times$)
    \item The standard errors also increase after shrinkage, which is counterintuitive
\end{itemize}

\subsection{Specific Examples: AMIGO3 vs Control}

The AMIGO3 perturbation shows much milder aberration:

\begin{table}[h]
\centering
\caption{Examples of PyDESeq2 aberrant shrinkage (AMIGO3 vs control)}
\label{tab:examples_amigo3}
\begin{tabular}{lcccccc}
\toprule
Gene & Raw LFC & Raw SE & Shrunk LFC & Shrunk SE & Inflation & Dispersion \\
\midrule
AIDA      & $-$0.286 & 0.217 & $-$0.463 & 0.279 & 1.62$\times$ & 3.24 \\
HIST1H1B  & $-$0.448 & 0.173 & $-$0.657 & 0.189 & 1.47$\times$ & 1.38 \\
LINC00910 & $-$0.309 & 0.186 & $-$0.430 & 0.227 & 1.39$\times$ & 1.77 \\
CDK13     & $-$0.302 & 0.174 & $-$0.397 & 0.212 & 1.32$\times$ & 1.21 \\
LBH       & $-$0.363 & 0.159 & $-$0.463 & 0.178 & 1.27$\times$ & 0.78 \\
\bottomrule
\end{tabular}
\end{table}

The AMIGO3 examples have much lower dispersions ($\alpha < 4$), explaining the milder aberration.

\section{Comparison with crispyx}

crispyx implements apeGLM shrinkage using Gaussian approximation (the ``stats'' method), which is robust to dispersion estimation differences.

\begin{table}[h]
\centering
\caption{Shrinkage behavior comparison: PyDESeq2 vs crispyx}
\label{tab:comparison}
\begin{tabular}{lcc}
\toprule
Metric & PyDESeq2 & crispyx \\
\midrule
Genes with $|$shrunk$|$ $>$ $|$raw$|$ & 21.5\% & \textbf{0.0\%} \\
Max inflation ratio & 13.74$\times$ & 1.00$\times$ \\
Median shrinkage ratio & 0.83 & 0.59 \\
Numerical warnings & Yes (overflow) & No \\
\bottomrule
\end{tabular}
\end{table}

\subsection{Why crispyx is More Robust}

crispyx uses the ``stats'' method for apeGLM shrinkage, which:
\begin{enumerate}
    \item Uses Gaussian approximation to the likelihood instead of re-fitting the NB likelihood
    \item Computes the posterior mode by solving:
    \begin{equation}
        \frac{\beta - \hat{\beta}_{\text{MLE}}}{\text{SE}^2} + \frac{2\beta}{s^2 + \beta^2} = 0
    \end{equation}
    where $s$ is the Cauchy prior scale
    \item Is independent of the stored dispersion values
    \item Guarantees $|\hat{\beta}_{\text{shrunk}}| \leq |\hat{\beta}_{\text{MLE}}|$ for $s > 0$
\end{enumerate}

\section{Root Cause Analysis}

The aberrant shrinkage occurs because:

\begin{enumerate}
    \item \textbf{Dispersion trend fitting fails}: PyDESeq2's parametric dispersion trend $\alpha(m) = \frac{a_0}{m} + a_\infty$ fails to converge
    
    \item \textbf{Fallback to mean-based trend}: A mean-based dispersion trend is used instead, leading to poorly estimated gene-wise dispersions
    
    \item \textbf{High dispersion estimates}: Some genes get extremely high dispersion estimates ($\alpha > 30$), especially for lowly-expressed genes
    
    \item \textbf{Flat NB likelihood}: With high dispersion, the NB likelihood becomes very flat, making the posterior mode sensitive to numerical noise
    
    \item \textbf{Optimization divergence}: The apeGLM optimization (using the ``full'' method with NB re-fitting) converges to incorrect values far from the MLE
\end{enumerate}

\section{Recommendations}

\subsection{For PyDESeq2 Users}
\begin{itemize}
    \item Check for the dispersion trend fitting warning
    \item If present, consider using crispyx or the original R DESeq2
    \item Validate shrinkage results by checking that $|$shrunk$|$ $\leq$ $|$raw$|$ for all genes
\end{itemize}

\subsection{For crispyx Users}
\begin{itemize}
    \item The default \texttt{method='stats'} is recommended for robustness
    \item Avoid \texttt{method='full'} unless you have validated your dispersion estimates
\end{itemize}

\section{Conclusion}

We have documented a significant numerical issue in PyDESeq2's apeGLM implementation that causes up to 21\% of genes to have \emph{inflated} rather than shrunk LFC estimates. This issue is triggered by dispersion trend fitting failure, which is common in single-cell CRISPR screen data. crispyx's implementation using Gaussian approximation is more robust and maintains the fundamental property that shrinkage should reduce, not increase, effect size magnitudes.

\begin{thebibliography}{9}
\bibitem{apeglm}
Zhu, A., Ibrahim, J.G., \& Love, M.I. (2019).
Heavy-tailed prior distributions for sequence count data: removing the noise and preserving large differences.
\textit{Bioinformatics}, 35(12), 2084-2092.
\end{thebibliography}

\end{document}
